Um die Nullkomponente des Stroms $j^0 = \rho$ für eine planare Wellenlösung der Klein-Gordon-Gleichung zu berechnen, betrachten wir die Lösung der Form $\phi(x) = A \exp(\mathrm{i} p x)$, wobei $p x = p_\mu x^\mu = p_0 t - \mathbf{p} \cdot \mathbf{x}$ und $A$ eine Konstante ist.

Die Klein-Gordon-Gleichung lautet:

\[
(\Box + m^2) \phi = 0
\]

wobei $\Box = \partial^\mu \partial_\mu = \frac{\partial^2}{\partial t^2} - \nabla^2$ der d'Alembert-Operator ist. Für die planare Wellenlösung ergibt sich:

\[
\Box \phi = -p_\mu p^\mu \phi = -(p_0^2 - \mathbf{p}^2) \phi
\]

Setzt man dies in die Klein-Gordon-Gleichung ein, erhält man die Dispersionsrelation:

\[
p_0^2 = \mathbf{p}^2 + m^2
\]

Der vierdimensionale Strom $j^\mu$ ist definiert als:

\[
j^\mu = \frac{\mathrm{i}}{2} \left( \phi^* \partial^\mu \phi - \phi \partial^\mu \phi^* \right)
\]

Für die Nullkomponente $j^0 = \rho$ ergibt sich:

\[
\rho = \frac{\mathrm{i}}{2} \left( \phi^* \partial^0 \phi - \phi \partial^0 \phi^* \right)
\]

Berechnen wir die zeitliche Ableitung:

\[
\partial^0 \phi = \frac{\partial \phi}{\partial t} = \mathrm{i} p_0 \phi
\]

und analog:

\[
\partial^0 \phi^* = -\mathrm{i} p_0 \phi^*
\]

Setzen wir diese Ausdrücke in die Gleichung für $\rho$ ein:

\[
\rho = \frac{\mathrm{i}}{2} \left( \phi^* (\mathrm{i} p_0 \phi) - \phi (-\mathrm{i} p_0 \phi^*) \right)
\]

\[
= \frac{\mathrm{i}}{2} \left( \mathrm{i} p_0 \phi^* \phi + \mathrm{i} p_0 \phi \phi^* \right)
\]

\[
= \frac{\mathrm{i}}{2} \cdot 2 \mathrm{i} p_0 \phi^* \phi
\]

\[
= p_0 |\phi|^2
\]

Da $|\phi|^2 = |A|^2$ konstant ist, folgt, dass $\rho = p_0 |A|^2$. Diese Nullkomponente des Stroms entspricht der Energiedichte der Welle, was im Einklang mit der Interpretation der Klein-Gordon-Gleichung als relativistische Wellengleichung steht.